\documentclass[11pt]{article}
\usepackage[T1]{fontenc}
\usepackage[cmex10]{amsmath}
\usepackage{amssymb}
\usepackage[margin=1in]{geometry}
\usepackage[hidelinks]{hyperref}
\DeclareFontShape{OMX}{cmex}{m}{n}{<->cmex10}{}
\hypersetup{pdftitle={Three-pole interlacing for ordered partitions and the largest zeros at arbitrary thresholds},pdfauthor={ProgramComputer}}
\title{Three-pole interlacing for ordered partitions and the largest zeros at arbitrary thresholds}
\author{ProgramComputer}
\date{November 25, 2026}
\begin{document}
\maketitle
\begin{abstract}
We prove that the polynomials of ordered set partitions counted by blocks of size at least three have simple negative zeros and consecutive rows strictly interlace in every size. Their generating function has three competing poles. A divided-difference representation of the remaining factor gives a quantitative gap in modulus. A linear combination of consecutive rows cancels the real pole, giving strict signs at the preceding roots. The analytic proof applies from row 94 and two exact coefficient algorithms certify the finite prefix. We determine the limiting zero distribution and the lower-edge scale. For each fixed block threshold, a reversed coefficient limit proves that every fixed number of outermost zeros is eventually simple, real and strictly interlacing, with explicit constants determined by exponential-series sections.
\end{abstract}
\section{Ordered set partitions counted by blocks of size at least three}

Let $O_n(t)=O_{n,3}(t)$ count ordered set partitions of $[n]$, with weight $t$ for each block of size at least three. The constant and leading coefficients are positive, and the degree is $\lfloor n/3\rfloor$. Alexandersson conjectures this property for every fixed threshold $j\ge2$~\cite{original}.

\textbf{Theorem.} For every $n\ge0$ all zeros of $O_n$ are simple and negative, and consecutive rows strictly interlace wherever there are roots to compare. In positive moduli, the $j$th root of row $n+1$ is smaller than the $j$th root of row $n$, which is smaller than the $(j+1)$st root of row $n+1$ whenever that root exists.

The analytic proof starts at $n=94$. The finite verification checks rows $0\le n\le94$ and consecutive pairs beginning at $0\le n\le93$.

\section{Poles}

Put $v=1/(1+q)$ for $t=-q<0$. The row exponential generating function is $1/((1+q)D_v(z))$, where

\[
D_v(z)=(1-v)e^z-1-z-z^2/2+2v.
\tag{1}
\]

Let $v_*$ be the unique positive zero in $(0,1)$ of $2\sqrt v-\log(1+2\sqrt v)+\log(1-v)$. In the variable $b=\sqrt v$ the expression starts at zero, tends to $-\infty$ at one, and its derivative has the sign of $1-2b-2b^2$ for $b>0$. Also $v_*<1/3$, since $23/20<2/\sqrt3<7/6$ and $e^{7/6}<129/40=(3/2)(1+23/20)$. The exponential bound follows by summing through degree seven and bounding the remaining tail by $(7/6)^8/(8!\,(1-7/54))$. We use $0<v<v_*$, or $q>q_*:=1/v_*-1$.

There is one real negative pole $-R$: $D_v(-r)$ starts at $v>0$, is strictly concave for $r\ge0$, and tends to $-\infty$. Its derivative in $r$ is negative at $R$. Also $R<9/5$ because $D_v(-9/5)$ increases with $v$ and at $v=1/3$ is $(2/3)e^{-9/5}-23/150<0$. This uses $e^{9/5}>1+9/5+(9/5)^2/2>100/23$. There is no positive real pole. Indeed, $D'_v(s)$ starts negative, decreases until $D_v''(s)$ changes sign, and then increases to infinity. Its positive zero $s_c$ satisfies $s_c<2\sqrt v$, since $D'_v(2\sqrt v)>0$ by the defining inequality for $v<v_*$. At the minimum $D_v(s_c)=2v-s_c^2/2>0$.

For an upper-half-plane pole $z=(u-1)+i\theta$, the imaginary equation is $(1-v)e^{u-1}\sin\theta=u\theta$. If $\sin\theta\le0$, then $\theta\ge\pi$, $u\le0$ and $|u|\le e^{-1}/\theta$. The real part $(u^2+1-\theta^2)/2-2v$ of the quadratic is then less than $-1$, while the exponential has modulus at most $e^{-1}$. Thus $u>0$ and $\theta\in(2\pi k,(2k+1)\pi)$ for some $k\ge0$.

Put $b=\theta\cot\theta$ and $d=\theta^2\csc^2\theta$. The real and imaginary equations together give

\[
u=b\pm\sqrt{d-1+4v},\qquad
G_v(\theta):=u-1-\log u-\log(\theta/\sin\theta)+\log(1-v)=0.
\tag{2}
\]

On each positive $u$ branch, differentiating $u^2-2bu+1-\theta^2-4v=0$ gives

\[
\theta(u-b)G'_v(\theta)=-u(b-1)^2-\theta^2(u-1)^2/u<0.
\tag{3}
\]

The right side can vanish only if $u=b=1$, which would contradict the quadratic. The plus branch is positive throughout $(0,\pi)$, and $G_v(0)>0$. At $\theta=1$, the alternating Taylor bounds give $u\ge2\cot1>1$, $\cot1\le13/20$, $\csc^2 1-1\le11/25$ and $\sin1\le101/120$. On $v\in[j/10,(j+1)/10]$, $j=0,1,2,3$, these give $u<U_j$, with $U_j=8/5,9/5,2,21/10$. Using $\log x\ge2(x-1)/(x+1)$ for $x\ge1$ and $\log(1-v)\le-v$,

\[
G_v(1)\le\frac{(U_j-1)^2}{U_j+1}-\frac{38}{221}-\frac j{10}<0.
\]

Therefore (3) gives exactly one plus-branch root with $0<\theta<1$. It has $u>1$, since the branch starts at $1+2\sqrt v$ and at any first crossing of $u=1$ its $G_v$ value is already negative. A positive minus branch in $(0,\pi)$ occurs only when $v<1/4$ and $0<\theta<\sqrt{1-4v}$. There $G_v$ increases by (3). Its initial value $-2\sqrt v-\log(1-2\sqrt v)+\log(1-v)$ is positive, since the derivative in $\sqrt v$ is positive on $(0,1/2)$. This excludes any pole on that branch. On each higher interval $k\ge1$, the product $1-\theta^2-4v$ of the quadratic roots in $u$ is negative. Hence only the plus branch is positive, and there is at most one pole by (3).

Write $z=s+i\theta$ for the unique principal upper pole, with $s>0$ and $0<\theta<1$. All complex poles are simple because at a zero of $D_v$,

\[
D'_v(z)=z^2/2-2v,
\tag{4}
\]

which can vanish only on the real axis. Conjugation supplies the lower poles. The other upper poles have imaginary parts greater than $2\pi k$, at most one for each $k\ge1$.

\section{Moduli and the negative real pole}

Factor $D_v(w)=(w-z)(w-\bar z)(w+R)g(w)$. The quadratic subtracted from $(1-v)e^w$ in (1) is its interpolating polynomial at $z,\bar z,-R$. The exponential interpolation remainder, in simplex integral form, gives

\[
g(w)=(1-v)\int_{\alpha+\beta+\gamma+\delta=1\atop \alpha,\beta,\gamma,\delta\ge0}
\exp(\alpha z+\beta\bar z-\gamma R+\delta w)\,d\sigma ,
\tag{5}
\]

where $d\sigma$ is simplex coordinate measure of total mass $1/6$. This is the Hermite--Genocchi formula for the third divided difference of the exponential, including repeated arguments by continuity. See de Boor~\cite[eqs. (50), (52)]{deboor}.

For real $w$, averaging under interchanging $\alpha,\beta$ replaces the phase by the positive factor $\cos((\alpha-\beta)\theta)$. Hence $c_1=(\log g)'(0)$ and $c_2=(\log g)''(0)$ are the mean and variance of $\delta\in[0,1]$ under a positive probability measure. In particular $c_1\ge0$ and $c_2\le c_1(1-c_1)$. Also $D_v(0)=v$ and $D'_v(0)=D''_v(0)=-v$. The logarithmic derivatives of the factorization give, with $r_0=|z|$ and $h=1+c_1$,

\[
\frac{2s}{r_0^2}=\frac1R+h,
\qquad
\frac{2}{r_0^2}=\frac1{R^2}+\left(\frac1R+h\right)^2-2-c_2 .
\]

Using the variance bound in the second identity gives

\[
\frac{R^2}{r_0^2}
=1+hR+\frac{h^2-2-c_2}{2}R^2
\ge 1+R-R^2/2
\ge 1+R/10 .
\tag{6}
\]

Thus $r_0<R$. Write $r=r_0/R$. As $R<9/5$ and $\log(1+R/10)\ge R/(10+R)>R/12$, we have $r\le e^{-R/24}$. Comparing (5) at $w=z$ and $w=-R$ gives

\[
\frac{|g(z)|}{g(-R)}
\le\frac{e^{s+R}}{\cos\theta}\le e^{3R}.
\tag{7}
\]

Indeed, the complex integral at $z$ costs at most $e^{s+R}$ relative to the underlying positive integral at $-R$, where the cosine is at least $\cos\theta$. Here $s<R$, $\theta<R$ and $-\log\cos\theta\le\theta^2<R$ by $\theta<1$, $\cos\theta\ge1-\theta^2/2$ and $-\log(1-x)\le x/(1-x)$.

We also have $\varphi:=\arg z<\pi/3$. For $v\to0$, the negative root $R$ tends to zero since $D_0(-r)<0$ for $r>0$ by the exponential Taylor remainder. By (6), $z\to0$, and (1) gives $z^3/v\to-6$, so the first-quadrant root has $\varphi\to\pi/3$. For all $v$ in the interval write $D_v(z)=v(2-e^z)+z^3 I(z)$ with $I(z)=\frac12\int_0^1(1-t)^2 e^{tz}\,dt$. The phase of $I(z)$ lies in $(0,\theta)$, and $2-e^z$ has negative imaginary part. Since $D_v(z)=0$, the continuous quantity $3\varphi+\arg I(z)-\arg(2-e^z)$ is an odd multiple of $\pi$. It tends to $\pi$ as $v\to0$, so equals $\pi$ throughout. Both additional terms are positive, giving $0<\varphi<\pi/3$.

\section{Uniform residue margins and sign alternation}

Put $N=n+1$ and normalize the principal conjugate residues by

\[
M_n(q)=\frac{2}{(1+q)|D'_v(z)|\,r_0^N},\qquad
\Phi_n(q)=N\varphi+\arg D'_v(z),\qquad
B=\frac{\sin\varphi}{|1+z/R|}.
\]

Here $\arg D'_v(z)\in(0,\pi)$ by (4) and $s>0$. The residue sum for $O_n(-q)/n!$ follows by integrating $1/((1+q)w^N D_v(w))$ over squares with real and imaginary coordinates $\pm 2\pi M$. On horizontal sides the imaginary part of $D_v(w)$ is $-(1+\Re w)\Im w$. Where $|1+\Re w|\ge1$ this bounds its modulus below by $2\pi M$. On the rest its real part grows quadratically in $M$. On vertical sides the exponential dominates on the right and the quadratic term dominates on the left. The integrals tend to zero, and the other complex residues are absolutely summable by (4) and their imaginary-part bounds. The contribution of every pole to the coefficient is minus its residue, since the residue at the origin equals the coefficient.

The normalized coefficient has the form $-\cos\Phi_n$ plus the normalized negative real residue and an error $E_N$ from the other complex poles. The negative real residue has absolute value

\[
\frac{|D'_v(z)|}{2D'_v(-R)}r^N
=B\frac{|g(z)|}{g(-R)}r^{N+1}\le B e^{-R}
\qquad(n\ge94).
\tag{8}
\]

The identity follows by differentiating the three-root factorization at $z,-R$. The bound uses $N+1\ge96$, (6), and (7). For (5) at $z$, $|g(z)|\le e^R/6<4/3$, and therefore $|D'_v(z)|\le(16/3)\theta R$. At every other upper pole $z_k$ of index $k\ge1$, (4) gives $|D'_v(z_k)|\ge17k^2$, since $|z_k|>2\pi k$, $v<1/3$ and $\pi>3$. Grouping complex conjugate residues and using $r_0<R<9/5$ then gives

\[
|E_N|\le \theta R(3/10)^N.
\tag{9}
\]

Here the coefficient factor $(16/51)\sum_{k\ge1}k^{-N-2}$ is less than one. Since $B\ge\theta/(2R)$, the total absolute normalized error in the cosine formula is bounded by

\[
B e^{-R}+|E_N|\le B(1-R/4)<1
\qquad(n\ge94).
\tag{10}
\]

Indeed, $1-e^{-R}\ge R/(1+R)>5R/14$ and $|E_N|/B\le(18/5)R(3/10)^N<R/16$. Also $B<1$.

As $q\downarrow q_*$, the principal $\theta$ tends to zero because $G_{v_*}(0)=0$ and (3) excludes any positive limiting angle. Formula (2) then gives $s=2\sqrt v+O(\theta^2)$, so $\arg D'_v(z)\to\pi/2$. At $q\to\infty$ we have $z\sim(6v)^{1/3}e^{i\pi/3}$ and $D'_v(z)\sim z^2/2$. Thus the phase runs continuously from $\pi/2$ to $(n+3)\pi/3$. For $m=\lfloor n/3\rfloor$, prepare ordered disjoint brackets with phase inside $(j\pi,(j+1)\pi)$, $j=1,\ldots,m$, using the last hit of the lower level before the first hit of the upper one. These levels occur in the endpoint range except for the last upper level when $n=3m$. If that level has no finite hit, use the last hit of $m\pi$ followed by the ray to infinity. Its leading sign at infinity is $(-1)^m$, opposite to the sign at $m\pi$. In all other brackets (10) gives opposite signs at the finite endpoints. Thus there is exactly one simple negative root in each bracket by the degree, and $\sin\Phi_n$ has sign $(-1)^j$ at the $j$th root in increasing modulus.

At such a root, (10) gives $|\cos\Phi_n|\le B(1-R/4)$. Add $1/R$ times the exponential coefficient $O_n(-q)/n!$ to $O_{n+1}(-q)/(n+1)!$. The negative real pole cancels exactly. After dividing by $M_{n+1}$, the principal complex terms become

\[
\sin\Phi_n\sin\varphi-(\cos\varphi+r)\cos\Phi_n,
\]

with error at most $2\theta R(3/10)^N$. Let $A=\sin\varphi$ and $C=\cos\varphi+r>1/2$, so that $B=A/\sqrt{A^2+C^2}$. The sine term has magnitude exceeding the cosine term by at least

\[
A\sqrt{1-B^2}-CB(1-R/4)=CBR/4\ge\theta/16.
\]

The remaining error is strictly smaller. Since the old coefficient vanishes at the tested point, the sign of $O_{n+1}(-q)$ is $(-1)^j$. From the positive value at $q=0$, these signs force one root before the first old root and one between each consecutive pair. If the degree increases, the leading sign at infinity forces one further root after the last old root. This proves strict consecutive interlacing for pairs beginning at $n\ge94$.

\section{Finite verification}

The finite verifier reconstructs the finite prefix by the first-block recurrence and by the delayed recurrence

\[
O_{n+1}=(n+1)O_n+2t O_n'+(n+\tbinom n2)O_{n-1}
 +\tbinom n2 t\bigl(O_{n-2}+(t-1)O'_{n-2}\bigr).
\]

The latter follows by extracting coefficients of $(1-z-z^2/2)H_z=(2+(t-1)z^2/2)tH_t+(1+z+t z^2/2)H$, where $H$ is the row generating function. Positive coefficients and exact Sturm counts give simple negative roots for all 95 rows. For each consecutive pair the signed Euclidean remainder degrees drop one at a time with positive leading coefficients and end at a nonzero constant. This verifies the orientation by evaluating the recursion at each root from the preceding step. At equal degrees one first replaces the old row by its remainder on division by the new row. The remainder has the same values at the new roots, so the sign condition puts one old root before the leftmost new root and one in each gap. The check also verifies the displayed rational comparisons. An independent check through $n=96$ used associated Stirling numbers for the large blocks and then added all ordered singleton and doubleton completions.

\section{Distribution of zeros at ordered threshold three}

Use the pole notation of the preceding sections. The critical size parameter is $v_*$, the positive double pole is $s_*=2\sqrt{v_*}$, and $q_*=v_*^{-1}-1$. For $q>q_*$, write $z(q)$ for the principal pole with positive imaginary part. Write $q_{n,h}$ for the $h$th smallest zero modulus.

\textbf{Corollary.} The empirical distribution functions of the zero moduli converge uniformly to $0$ for $q\le q_*$ and to $3\arg z(q)/\pi$ for $q>q_*$. For fixed $h\ge1$,
\[
\lim_{n\to\infty} n^2(q_{n,h}-q_*)
 =\frac{s_*^3}{2(e^{s_*}-2)v_*^2}\pi^2 h^2 .
\]
The complementary $n^3$ scale at the outer edge is given by the exponential-section limit in the next section.

\textbf{Proof.} For $n\ge94$ the threshold-three proof gives a root in each phase band $(h\pi,(h+1)\pi)$, exhausting the degree. At any $q>q_*$ the bands below the current phase have been crossed before $q$, and the bands above it must be crossed after $q$. At the last band in degree $n=3m$ the leading sign at infinity supplies the endpoint whenever the finite upper level is absent. Degree exhaustion and the phase endpoint bounds give, uniformly,
\[
\#\{h:q_{n,h}\le q\}=(n+1)\arg z(q)/\pi+O(1).
\]
All roots in those rows exceed $q_*$. This proves uniform convergence. The distribution is strictly increasing on its support, by analyticity of the principal branch, monotonicity inherited from the counts, and the different endpoint values.

At fixed $h$ the band gives $\arg z(q_{n,h})=O(1/n)$, hence $q_{n,h}\to q_*$. Near $q_*$ the ratio of the principal modulus to the negative-real pole modulus is bounded strictly below one by (6) in the threshold-three proof, as $R$ tends to a positive value. The contributions of that real pole and the distant complex poles therefore tend to zero uniformly there. The root phase is consequently $(h+1/2)\pi+o(1)$. As $\arg D'_v(z)\to\pi/2$, we get $(n+1)\arg z\to h\pi$.

At the positive double pole, $D_{zz}=s_*>0$ and $\partial_v D=2-e^{s_*}<0$. The latter sign also follows from $e^{s_*}=(1+s_*)/(1-s_*^2/4)$ and $s_*>1$: the defining function $2\sqrt v-\log(1+2\sqrt v)+\log(1-v)$ is positive at $v=1/4$ because $e>8/3$. For $v<v_*$ near $v_*$ the two nonreal poles split from this double zero. Since $\partial_vD\ne0$, one can solve analytically for $v=v(z)$ there. Its Taylor series is
\[
v-v_* = \frac{s_*}{2(e^{s_*}-2)}(z-s_*)^2+O((z-s_*)^3).
\]
The two square-root branches thus have $\theta=\Im z\sim\sqrt{2(e^{s_*}-2)(v_*-v)/s_*}$ and $\arg z\sim\theta/s_*$. Finally $v_*-v\sim v_*^2(q-q_*)$, yielding the asserted limit. $\square$


\section{The largest zeros at every fixed ordered threshold}

Fix $j\ge2$. Write $A_{n,k}=[t^k]O_{n,j}(t)$ and $n=jm+s$, $0\le s<j$. Set
\[
\alpha=1+\frac1{j+1},\quad L_m=j!\,(\alpha m)^j,\qquad
 H_s(x)=\sum_{h\ge0}\frac{x^h}{(jh+s)!}.
\]
Let $0<\lambda_{s,1}<\lambda_{s,2}<\cdots$ be the moduli of the negative zeros of $H_s$.

\textbf{Theorem.} For every finite $K$, the $K$ largest-modulus complex zeros of $O_{n,j}$ are simple and negative for all sufficiently large $n$. Their moduli $Q_{n,1}>\cdots>Q_{n,K}$ satisfy $Q_{n,i}/L_m\to1/\lambda_{s,i}$ along each residue $s$. They strictly alternate with those in the next row at this end. With indices within the fixed range, the order for $s<j-1$ is
\[
Q_{n,i}>Q_{n+1,i}>Q_{n,i+1},
\]
and for $s=j-1$ it is
\[
Q_{n+1,i}>Q_{n,i}>Q_{n+1,i+1}.
\]
This conclusion allows every fixed threshold and unbounded row size. It describes the outermost complex zeros even if real-rootedness of the remaining polynomial is unknown.

\textbf{Proof.} Write $S(z)=\sum_{a=1}^{j-1}z^a/a!$ and $C(z)=\sum_{a\ge j}z^a/a!$. Let $T=1/(1-S)$ and $W=j!z^{-j}C T$, analytic at zero. The ordered-block generating function gives
\[
A_{n,k}/n!=(j!)^{-k}[z^{n-jk}]\,T(z)W(z)^k .
\]
This coefficient extraction also appears in the public computation~\cite{prior}. The functions $T$ and $W$ have nonnegative coefficients, constant term one, and $W'(0)=\alpha$. For fixed $h,e\ge0$,
\[
m^{-e}[z^e]T(z)W(z)^{m-h}\longrightarrow \alpha^e/e!.
\]
Indeed $T(w/m)W(w/m)^{m-h}\to e^{\alpha w}$ locally uniformly in $w$ since $\log W(z)=\alpha z+O(z^2)$ near zero. Applying coefficient convergence at $e=s+jh$ and at $e=s$ gives
\[
\frac{A_{n,m-h}}{A_{n,m}L_m^h}\longrightarrow \frac{s!}{(s+jh)!}.
\]

The convergence holds for the whole reversed polynomial locally uniformly. For any fixed $B>0$, the nonnegative coefficients and evaluation at $z=B/m$ bound
\[
m^{-e}[z^e]T(z)W(z)^{m-h}
 \le B^{-e}T(B/m)W(B/m)^m\le C_B B^{-e}
\]
for $0\le h\le m$ and sufficiently large $m$. Also $m^{-s}[z^s]T(z)W(z)^m$ stays bounded below by a positive constant. Taking $B$ arbitrarily large therefore gives a summable geometric majorant on each compact disk for the terms of
\[
\mathcal R_{m,s}(x)=\sum_{h=0}^m\frac{A_{n,m-h}}{A_{n,m}L_m^h}x^h
 \ \longrightarrow\ s!H_s(x).
\]
These are the exponential-series sections. The Gauss multiplication formula writes $s!H_s$ as a hypergeometric ${}_0F_{j-1}$ at $x/j^j$, with the positive parameters $(s+1)/j,\ldots,(s+j)/j$ and the single parameter $1$ removed. The hyper-Bessel zero theorem puts all its zeros on the negative real axis. The series has order at most $1/j<1$ and is nonpolynomial, so Hadamard factorization gives infinitely many zeros and a genus-zero product. See Baricz and Singh~\cite[Theorem 2]{baricz}.

Here is a short proof of the strict ordering needed in the transfer. For $s>0$, $H_{s-1}=(jx\partial_x+s)H_s$, while $H_{j-1}=jH_0'$. At a negative zero of multiplicity $b$ either operation lowers its multiplicity by one. At points $x<0$ where the input is nonzero, a newly created zero is simple, since
\[
\left(\frac{H_s'}{H_s}+\frac{s}{jx}\right)'=-\sum_{H_s(y)=0}\frac{1}{(x-y)^2}-\frac{s}{jx^2}<0 .
\]
For the derivative step take $s=0$. Traversing the whole cycle returns to the original function, contradicting an initial multiplicity greater than one: it either decreases throughout, or drops to at most one and can never exceed one again. Thus all zeros are simple. In each gap between successive input zeros the displayed logarithmic expression decreases from positive to negative infinity. In the rightmost gap ending at zero it has a unique zero for $s>0$ and none for $s=0$, by the genus-zero product. These relations give $\lambda_{s-1,i}<\lambda_{s,i}$ and $\lambda_{j-1,i}<\lambda_{0,i+1}$, hence $\lambda_{0,i}<\cdots<\lambda_{j-1,i}<\lambda_{0,i+1}$.

Choose a large disk containing the first $K$ limit zeros and no further zero, then disjoint smaller disks centered on them. Rouché's theorem on all these boundaries gives one simple zero in each small disk and exactly $K$ zeros in the larger one for all large $m$. Conjugation gives reality. Reversing takes a zero $x$ to $L_m/x$. The strict inequalities between section zeros and $L_{m+1}/L_m\to1$ give both consecutive orders, including the degree increase at the residue wrap. There are only $j$ residues. At $j=2$ these constants agree with the two explicit parity formulas for nonsingleton blocks~\cite{companion}. $\square$

\begin{thebibliography}{9}
\bibitem{original} P. Alexandersson, \emph{Real-rooted polynomials: open problems}, The symmetric functions catalog, section on set partitions and trees, \href{https://www.symmetricfunctions.com/realRootedOpenProblems.htm\#openProblemsSetPartitions}{author's conjecture}, accessed November 25, 2026.
\bibitem{deboor} C. de Boor, Divided differences, \emph{Surveys in Approximation Theory} \textbf{1} (2005), 46--69. \href{https://pages.cs.wisc.edu/~deboor/sat/papers/2/2.pdf}{Author's copy}.
\bibitem{baricz} \'{A}. Baricz and S. Singh, Zeros of some special entire functions, \href{https://arxiv.org/abs/1702.00626}{arXiv:1702.00626} (2017).
\bibitem{prior} \emph{Alexandersson's ordered-set-partition conjecture}, public investigation, \href{https://github.com/1337isnot1337/solved-conjectures/issues/108}{issue 108 and linked research note}.
\bibitem{companion} ProgramComputer, \emph{Interlacing and limiting zeros of ordered set partitions counted by nonsingleton blocks}, \href{https://github.com/ProgramComputer/open-math-research/blob/main/topics/ordered-large-blocks/paper.pdf}{companion manuscript} (2026).
\end{thebibliography}
\end{document}
